\chapter{Geltungsbereich und Grundbegriffe}

Dieser Leitfaden ergänzt Band~5 um eine ausgewählte Übersicht
wichtiger Definitionen, Charakterisierungen und Folgerungen für
\emph{allgemeine Funktionen}. Seine Grundlagen sind die Mengenlehre
in ZF, die Totalität aus Band~4 und die funktionale Eindeutigkeit
aus Band~5. Die Aussagen benötigen kein Auswahlaxiom.
Beweise und Beweisskizzen entfallen. Das Format folgt den Leitfäden
zu Band~3, Band~4, Band~28 und Band~37.

\section{Gegenstand und Abgrenzung}

Erfaßt werden die Grundbegriffe und die wichtigsten Regeln für
Funktionsgraphen, Bilder, Urbilder, Einschränkungen, Zusammenfügen,
Komposition und mengenbasierte Funktionskonstruktionen.
Die Ergebnisse gelten für beliebige Mengen; Endlichkeit ist
keine allgemeine Voraussetzung. Eine vollständige Aufzählung aller
bekannten oder ableitbaren Aussagen wird nicht beansprucht.

Injektive, surjektive und bijektive Funktionen sowie ihre besonderen
Charakterisierungen bleiben den Bänden~6--8 vorbehalten.
Ordnungs- und Äquivalenzrelationen, Auswahlprinzipien, Zahlenfolgen,
algebraische und topologische Strukturen werden hier nicht entwickelt.
Auch Identitäts-, Inklusions-, Projektions- und Umkehrfunktionen
bleiben gemäß der Gliederung des Hauptbands in ihren eigenen Bänden.

Eine Bedingung innerhalb eines Satzes, etwa $C\subseteq A$ oder
die Übereinstimmung zweier Funktionen auf einer Überlappung,
ist eine ausdrücklich genannte Voraussetzung dieses Satzes.
Sie ist kein weiteres Grundaxiom der Funktionstheorie.
Allgemeine Rechenregeln werden aufgenommen, soweit sie unmittelbar
zum Arbeiten mit beliebigen Funktionen gehören.

\section{Funktion, Wert und Graph}

\begin{definition}[Funktion als totale eindeutige Relation]
\label{b05guide:def:function}
Für Mengen $A,B,f$ bedeutet $f\colon A\to B$, daß
\[
 \begin{gathered}
 f\subseteq A\times B,\\
 \forall x\in A\,\exists y\in B\,((x,y)\in f),\\
 \forall x\in A\,\forall y,z\in B\,
   \bigl((x,y)\in f\land(x,z)\in f\Rightarrow y=z\bigr).
 \end{gathered}
\]
$A$ heißt \emph{Definitionsbereich}, $B$ \emph{Zielmenge}.
Die zweite Bedingung ist die Totalität, die dritte die
funktionale Eindeutigkeit.
\end{definition}

\begin{definition}[Funktionswert und Graphnotation]
Für $f\colon A\to B$ und $x\in A$ bezeichnet $f(x)$ das
eindeutig bestimmte $y\in B$ mit $(x,y)\in f$.
Außerhalb von $A$ wird $f(x)$ nicht verwendet.
Die Schreibweise $x\mapsto t(x)$ beschreibt eine Funktion durch
ihre Werte, sofern für jedes $x$ des angegebenen Definitionsbereichs
ein eindeutiger Wert in der angegebenen Zielmenge feststeht.
\end{definition}

\begin{theorem}[Auswertung und Rekonstruktion des Graphen]
\label{b05guide:thm:graph}
Für $f\colon A\to B$, $x\in A$ und $y\in B$ gilt
\[
 f(x)\in B,\qquad (x,y)\in f\ \Longleftrightarrow\ y=f(x).
\]
Insbesondere ist
\[
 f=\{(x,f(x))\mid x\in A\}.
\]
\end{theorem}

Wie im Hauptband ist die Funktion selbst eine Menge geordneter
Paare. Die Zielmenge ist kein zusätzlicher Bestandteil ihres
Graphen. Deshalb kann derselbe Graph mehrere zulässige Zielmengen
haben. Diese Konvention ist bei Gleichheitsaussagen maßgeblich.

\section{Orientierung}

Kapitel~\ref{b05guide:ch:criteria} charakterisiert Funktionen
und klärt ihre Existenz und Gleichheit.
Kapitel~\ref{b05guide:ch:images} sammelt Bild-, Urbild- und Faserregeln.
Kapitel~\ref{b05guide:ch:restrictions} behandelt Einschränkungen,
Zusammenfügen und Erweiterungen.
Kapitel~\ref{b05guide:ch:composition} beschreibt Komposition
und paarwertige Funktionen.
Kapitel~\ref{b05guide:ch:functionsets} behandelt Funktionenmengen
und mengenwertige Funktionen.

Die Aussage-Nummern mit Präfix $5H$ gehören ausschließlich zu
diesem Leitfaden. Sie verändern die Nummerierung des Fachbands nicht.
